\documentclass[11pt]{article}
\usepackage{amsmath,amssymb,amsthm}
\newtheorem{theorem}{Theorem}[section]
\newtheorem{lemma}[theorem]{Lemma}
\newtheorem{proposition}[theorem]{Proposition}
\newtheorem{corollary}[theorem]{Corollary}
\theoremstyle{definition}
\newtheorem{definition}[theorem]{Definition}

\title{Triangle Numbers, Second Version}
\author{A synthetic example}
\date{}

\begin{document}
\maketitle

\section{Definitions}

\begin{definition}[Triangle numbers]\label{def:triangle}
For a natural number $n$, the triangle number $T(n)$ is the sum $0 + 1 + \dots + n$.
\end{definition}

\section{Two facts added in this version}

\begin{theorem}\label{thm:max}
For all natural numbers $m$ and $n$, $m + n \le 2 \max(m, n)$.
\end{theorem}

\begin{lemma}\label{lem:sq}
Every natural number $n$ satisfies $n \le n^2$.
\end{lemma}

\section{Bounds}

\begin{theorem}\label{thm:bound}
There is a constant $C$ such that for every natural number $n$, the triangle number satisfies
$T(n) \le C n^2$.
\end{theorem}

\begin{lemma}\label{lem:twice}
For every natural number $n$, twice the triangle number $T(n)$ equals $n(n+1)$.
\end{lemma}

\begin{proposition}\label{prop:succ}
For every natural number $n$, the triangle number $T(n+1)$ is $T(n) + (n + 1)$.
\end{proposition}

\begin{corollary}\label{cor:const}
For every natural number $n$ there is a constant $C$ such that the triangle number
satisfies $T(n) \le C$.
\end{corollary}

\end{document}
